% Options for packages loaded elsewhere
% Options for packages loaded elsewhere
\PassOptionsToPackage{unicode}{hyperref}
\PassOptionsToPackage{hyphens}{url}
%
\documentclass[
  ignorenonframetext,
  compress, aspectratio=149]{beamer}
\newif\ifbibliography
\usepackage{pgfpages}
\setbeamertemplate{caption}[numbered]
\setbeamertemplate{caption label separator}{: }
\setbeamercolor{caption name}{fg=normal text.fg}
\beamertemplatenavigationsymbolshorizontal
% remove section numbering
\setbeamertemplate{part page}{
  \centering
  \begin{beamercolorbox}[sep=16pt,center]{part title}
    \usebeamerfont{part title}\insertpart\par
  \end{beamercolorbox}
}
\setbeamertemplate{section page}{
  \centering
  \begin{beamercolorbox}[sep=12pt,center]{section title}
    \usebeamerfont{section title}\insertsection\par
  \end{beamercolorbox}
}
\setbeamertemplate{subsection page}{
  \centering
  \begin{beamercolorbox}[sep=8pt,center]{subsection title}
    \usebeamerfont{subsection title}\insertsubsection\par
  \end{beamercolorbox}
}
% Prevent slide breaks in the middle of a paragraph
\widowpenalties 1 10000
\raggedbottom
\AtBeginPart{
  \frame{\partpage}
}
\AtBeginSection{
  \ifbibliography
  \else
    \frame{\sectionpage}
  \fi
}
\AtBeginSubsection{
  \frame{\subsectionpage}
}
\usepackage{iftex}
\ifPDFTeX
  \usepackage[T1]{fontenc}
  \usepackage[utf8]{inputenc}
  \usepackage{textcomp} % provide euro and other symbols
\else % if luatex or xetex
  \usepackage{unicode-math} % this also loads fontspec
  \defaultfontfeatures{Scale=MatchLowercase}
  \defaultfontfeatures[\rmfamily]{Ligatures=TeX,Scale=1}
\fi
\usepackage{lmodern}

\usetheme[]{Montpellier}
\usecolortheme[]{dove}
\useinnertheme[]{rounded}
\useoutertheme[]{miniframes}
\ifPDFTeX\else
  % xetex/luatex font selection
\fi
% Use upquote if available, for straight quotes in verbatim environments
\IfFileExists{upquote.sty}{\usepackage{upquote}}{}
\IfFileExists{microtype.sty}{% use microtype if available
  \usepackage[]{microtype}
  \UseMicrotypeSet[protrusion]{basicmath} % disable protrusion for tt fonts
}{}
\makeatletter
\@ifundefined{KOMAClassName}{% if non-KOMA class
  \IfFileExists{parskip.sty}{%
    \usepackage{parskip}
  }{% else
    \setlength{\parindent}{0pt}
    \setlength{\parskip}{6pt plus 2pt minus 1pt}}
}{% if KOMA class
  \KOMAoptions{parskip=half}}
\makeatother


\usepackage{longtable,booktabs,array}
\usepackage{calc} % for calculating minipage widths
\usepackage{caption}
% Make caption package work with longtable
\makeatletter
\def\fnum@table{\tablename~\thetable}
\makeatother
\usepackage{graphicx}
\makeatletter
\newsavebox\pandoc@box
\newcommand*\pandocbounded[1]{% scales image to fit in text height/width
  \sbox\pandoc@box{#1}%
  \Gscale@div\@tempa{\textheight}{\dimexpr\ht\pandoc@box+\dp\pandoc@box\relax}%
  \Gscale@div\@tempb{\linewidth}{\wd\pandoc@box}%
  \ifdim\@tempb\p@<\@tempa\p@\let\@tempa\@tempb\fi% select the smaller of both
  \ifdim\@tempa\p@<\p@\scalebox{\@tempa}{\usebox\pandoc@box}%
  \else\usebox{\pandoc@box}%
  \fi%
}
% Set default figure placement to htbp
\def\fps@figure{htbp}
\makeatother





\setlength{\emergencystretch}{3em} % prevent overfull lines

\providecommand{\tightlist}{%
  \setlength{\itemsep}{0pt}\setlength{\parskip}{0pt}}



 


\usepackage{graphicx}
\usepackage{xcolor}
\usepackage{fancyvrb}
\usepackage{listings}
\definecolor{ColPrimario}{RGB}{100,8,58}
\definecolor{ColSecondario}{RGB}{20,143,184}
\setbeamercolor{structure}{fg=ColPrimario}
\setbeamercolor{itemize item}{fg=ColPrimario}
\setbeamercolor{itemize subitem}{fg=ColSecondario}
\setbeamercolor{palette primary}{bg=ColPrimario,fg=white}
\setbeamercolor{palette secondary}{bg=ColSecondario,fg=white}
\setbeamercolor{section in head/foot}{bg=ColPrimario,fg=white}
\setbeamercolor{subsection in head/foot}{bg=ColSecondario,fg=white}
\setbeamercolor{separation line}{fg=white}
\setbeamercolor{upper separation line head}{bg=white}
\setbeamercolor{middle separation line head}{bg=white}
\setbeamercolor{lower separation line head}{bg=white}
\lstset{basicstyle=\footnotesize\ttfamily, breaklines=true}
\DefineVerbatimEnvironment{Highlighting}{Verbatim}{commandchars=\\\{\}, fontsize=\footnotesize}
\DefineVerbatimEnvironment{verbatim}{Verbatim}{fontsize=\footnotesize}
\AtBeginSubsection{}
\AtBeginDocument{\institute[]{Università di Padova} }
\\AtBeginDocument{\\author{Ottavia M. Epifania\\\\\\texttt{ottavia.epifania@unipd.it} \\\\ Margherita Calderan\\\\\\texttt{margherita.calderan@unipd.it}}}
\AtBeginSection[]
        {
           \begin{frame}[plain]
           \tableofcontents[currentsection, hideallsubsections]
            \end{frame}
        }
\makeatletter
\@ifpackageloaded{caption}{}{\usepackage{caption}}
\AtBeginDocument{%
\ifdefined\contentsname
  \renewcommand*\contentsname{Table of contents}
\else
  \newcommand\contentsname{Table of contents}
\fi
\ifdefined\listfigurename
  \renewcommand*\listfigurename{List of Figures}
\else
  \newcommand\listfigurename{List of Figures}
\fi
\ifdefined\listtablename
  \renewcommand*\listtablename{List of Tables}
\else
  \newcommand\listtablename{List of Tables}
\fi
\ifdefined\figurename
  \renewcommand*\figurename{Figure}
\else
  \newcommand\figurename{Figure}
\fi
\ifdefined\tablename
  \renewcommand*\tablename{Table}
\else
  \newcommand\tablename{Table}
\fi
}
\@ifpackageloaded{float}{}{\usepackage{float}}
\floatstyle{ruled}
\@ifundefined{c@chapter}{\newfloat{codelisting}{h}{lop}}{\newfloat{codelisting}{h}{lop}[chapter]}
\floatname{codelisting}{Listing}
\newcommand*\listoflistings{\listof{codelisting}{List of Listings}}
\makeatother
\makeatletter
\makeatother
\makeatletter
\@ifpackageloaded{caption}{}{\usepackage{caption}}
\@ifpackageloaded{subcaption}{}{\usepackage{subcaption}}
\makeatother

\usepackage{bookmark}
\IfFileExists{xurl.sty}{\usepackage{xurl}}{} % add URL line breaks if available
\urlstyle{same}
\hypersetup{
  pdftitle={Probabilità, distribuzioni di probabilità, la distribuzione Normale e i test psicologici},
  pdfauthor={Ottavia M. Epifania, Margherita Calderan},
  hidelinks,
  pdfcreator={LaTeX via pandoc}}


\title{Probabilità, distribuzioni di probabilità, la distribuzione
Normale e i test psicologici}
\subtitle{Test per le organizzazioni}
\author{Ottavia M. Epifania, Margherita Calderan}
\date{}

\begin{document}
\frame{\titlepage}


\section{La probabilità}\label{la-probabilituxe0}

\subsection{Introduzione}\label{introduzione}

\begin{frame}{Introduzione}
\begin{itemize}
\item
  In termini generali, la \textbf{probabilità} può essere vista come il
  grado di \textbf{fiducia} rispetto al manifestarsi di un evento
\item
  In altre parole, la probabilità è una misura matematica che serve ad
  esprimere la nostra \textbf{incertezza}
\end{itemize}
\end{frame}

\subsection{Definizioni: Spazio campionario, esito e
evento}\label{definizioni-spazio-campionario-esito-e-evento}

\begin{frame}{Definizioni: Spazio campionario, esito e evento}
\begin{itemize}
\item
  Dato un fenomeno casuale (ad esempio, un lancio di un dado), l'insieme
  di tutti i possibili esiti (risultati) è detto \textbf{spazio
  campionario} (S)
\item
  Gli elementi di \textbf{S} devono essere tra loro \textbf{mutuamente
  esclusivi} e complessivamente \textbf{esaustivi}
\end{itemize}
\end{frame}

\begin{frame}
L'intersezione tra gli eventi è nulla, quindi se abbiamo due eventi A e
B se la loro intersezione è nulla questi eventi si dicono
\textbf{mutuamente esclusivi} o \textbf{mutuamente disgiunti}

\[A \cap B = \emptyset\]

Se l'unione tra gli eventi coincide con lo spazio campione, gli eventi
si dicono \textbf{collettivamente esaustivi}

\[(A + B) = S\]

\textbf{Partizione:} Se gli eventi sono sia mutuamente esclusivi che
collettivamente esaustivi, si può dire che ogni evento costituisce una
\textbf{partizione unica dello spazio campione}.
\end{frame}

\begin{frame}{Evento}
\phantomsection\label{evento}
Un \textbf{evento}(\textbf{E}) è un sottoinsieme di S:

\begin{itemize}
\tightlist
\item
  Fenomeno casuale: Un lancio di un dado a 6 facce
\item
  Spazio campionario: \(S =  \{1; 2 ; 3 ; 4 ; 5 ; 6 \}\)
\item
  Evento: L'uscita di un numero pari, \(E=\{2 ,4 ,6\}\)
\end{itemize}
\end{frame}

\subsection{Spazi campionari discreti e
continui}\label{spazi-campionari-discreti-e-continui}

\begin{frame}{Spazi campionari discreti e continui}
Uno spazio campionario si dice \textbf{discreto} se i suoi elementi sono
numerabili in categorie

\begin{itemize}
\tightlist
\item
  Fenomeno: Numero di errori in un compito
\item
  Spazio Campionario: \(S =  \{1; 2 ; 3 ; 4 ; 5 ; \ldots\}\)
\end{itemize}
\end{frame}

\begin{frame}
Uno spazio campionario si dice \textbf{continuo} se i suoi elementi non
sono direttamente numerabili, ma rappresentano un continuo di valori

\begin{itemize}
\tightlist
\item
  Fenomeno: tempo di risposta (\(\theta\))
\item
  Spazio Campionario: \(S = \theta \geq 0\)
\end{itemize}
\end{frame}

\subsection{La probabilità}\label{la-probabilituxe0-1}

\begin{frame}{La probabilità}
La probabilità, \(P()\), è una funzione matematica che assegna dei
valori numerici a degli esiti possibili di un fenomeno casuale

\begin{itemize}
\tightlist
\item
  Un valore di probabilità deve essere non negativo: \(P(A) \ge 0\).
\item
  La probabilità dell'intero spazio campionario è pari a 1:
  \(P(\Omega)=1\) (i.e., La somma delle probabilità di tutti gli esiti
  dello spazio campionario deve essere pari a 1).
\item
  Se due eventi sono \textbf{mutuamente esclusivi} (non possono
  verificarsi insieme), allora la probabilità che si verifichi l'uno
  \emph{oppure} l'altro è pari alla somma delle probabilità dei singoli
  eventi: se \(A \cap B = \varnothing\), allora
  \(P(A \cup B)=P(A)+P(B)\).
\end{itemize}

(Assiomi di Kolmogorv, 1956)
\end{frame}

\section{Distribuzioni di
probabilità}\label{distribuzioni-di-probabilituxe0}

\begin{frame}{Distribuzioni di probabilità}
Una distribuzione di probabilità è l'insieme di tutti gli esiti dello
spazio campionario e delle corrispettive probabilità.

\begin{itemize}
\item
  \textbf{Fenomeno}: \(X\) = esito del lancio di una moneta non
  truccata; dove \(X\) è una variabile aleatoria (casuale) che può
  assumere i valori \(x = 0\) (Croce) e \(x = 1\) (Testa).
\item
  \textbf{Spazio campionario}: \(S = \{0, 1\}\).
\item
  \textbf{Distribuzione di probabilità}: \[
  \Pr(X = x) = \theta^{x}(1-\theta)^{1-x},
  \] dove \(\theta\), il parametro che definisce la distribuzione di
  probabilità, è pari a \(0.5\).
\item
  \textbf{Proprietà della distribuzione di probabilità}: \[
  \Pr(X = 0) = 0.5 \ge 0;\quad \Pr(X = 1) = 0.5 \ge 0
  \] \[
  \Pr(X = 0) + \Pr(X = 1) = 1
  \]
\end{itemize}
\end{frame}

\section{La distribuzione normale}\label{la-distribuzione-normale}

\section{Credits}\label{credits}

\begin{frame}{Credits}
Altoè, G. (2022). Corso di Testing Psicologico, Scienze psicologiche
dello sviluppo, della personalità e delle relazioni interpersonali, A.A.
2022/23
\end{frame}




\end{document}
